\section{工程即方法：一致性、分块与可复现性}\label{sec:engineering}

本节不是“系统实现附录”。我们主张：在细长结构分割任务中，若干被当作工程细节的变量会实质改变结论；把它们提升为方法学对象、并用契约与协议固定下来，本身就是可复现研究的一部分。以下五节分别处理一致性、分块、显存、流水线与评测纪律，最后给出分层架构。

\subsection{患者级划分与泄漏审计}\label{sec:leakage}

MGD-1K 的同一名受试者常有多幅影像（上/下睑、左/右眼），文件名中的数字研究编号提供了可执行的患者标识。我们把整组受试者而不是单幅影像分配到 train/val/test，得到 \nTrain/\nVal/\nTest{} 图像（\nTrainPatients/\nValPatients/\nTestPatients{} 名患者），三部分患者互斥；分层依据是每组的最大专家 Meiboscore，以避免重度病例在某一 split 中缺失。划分文件同时记录文件清单签名 \splitSig{} 与实际派分签名 \assignmentSig{}：前者只回答数据集是否变化，后者回答本次实验到底用了哪一份派分。缺少后者时，同名文件可以被重新派分而无法追溯，这是“同名不同分”的根源。

泄漏审计给出三层证据。第一，官方 \officialTestImages{} 例 Meiboscore 测试清单来自 \officialTestPatients{} 名患者；其余 800 例中有 \officialLeakImages{} 例来自这些患者。第二，即使采用患者级划分，\nTrain{} 训练集中仍有 \officialLeakPatients/\officialTestPatients{} 名官方测试患者（355 张图像）。第三，构造 clean-official 子集，只保留官方测试中患者不在训练集的 \cleanOfficialImages{} 图/\cleanOfficialPatients{} 患者，并训练一个 image-level leaky 对照（\leakyTrainImages{} 训练 / \leakyValImages{} 验证）。三者在同一 clean-official 子集上的差，就是泄漏的代价。图~\ref{fig:pairing} 给出配对、患者分组与两种签名的关系。

\subsection{训练—推理一致性契约}\label{sec:contract}

深度分割模型的性能差异常被归因于网络，但同一权重在不同预处理/分块/阈值组合下可以给出截然不同的 Dice。为排除这一类隐性变量，本文把“训练看到什么”与“推理看到什么”写成可执行的契约，而不是文档约定。

契约由三条硬规则组成。

\begin{enumerate}[label=\textbf{C\arabic*.}]
\item \textbf{预处理逐像素一致。}训练时的在线变换与推理时的预处理调用同一个实现：灰度读取、EXIF 方向校正、双线性缩放到 $640\times1280$、除以 255、CLAHE（clip 0.02）都是共享函数；训练侧只在共享函数之后追加随机增强。单元测试对一组固定输入逐像素比对训练路径与推理路径，容差 $10^{-6}$，任何一侧改动都会使测试失败。
\item \textbf{分块参数统一。}训练分块、评估分块与部署默认值来自同一份 schema（默认 $384\times768$、重叠 0.25）。评估脚本不得自行改写分块；若实验需要改变分块，必须在报告中显式记录，并视为一次独立实验（见第~\ref{sec:tiling}~节）。
\item \textbf{三方数据配对硬失败。}影像、腺体掩膜、睑缘掩膜按 basename 建立索引；任一侧缺失、出现孤儿掩膜、重名或命名异常时，数据加载直接抛出完整性错误并终止，而不是跳过样本继续训练。静默跳过会让不同实验的“有效样本数”不同，从而使比较失去意义。
\end{enumerate}

图~\ref{fig:preprocess} 与图~\ref{fig:pairing} 分别给出一致性契约与三向配对/划分流程。契约的收益不是“代码更整洁”，而是让第~\ref{sec:results}~节的所有模型比较建立在同一输入分布上；任何不满足契约的运行都会在产物中留下可检查的痕迹。

\begin{figure}[t]
\centering
\resizebox{\linewidth}{!}{%
\begin{tikzpicture}[
  d/.style={draw=slate, fill=mist, rounded corners=1.2pt, align=center, font=\scriptsize, minimum height=9mm, minimum width=20mm, inner sep=2pt},
  ok/.style={draw=teal, fill=teal!12, rounded corners=1.2pt, align=center, font=\scriptsize, minimum height=9mm, minimum width=20mm},
  bad/.style={draw=alert, fill=alert!8, rounded corners=1.2pt, align=center, font=\scriptsize, minimum height=9mm, minimum width=20mm},
  arr/.style={-{Stealth[length=1.8mm]}, draw=slate, thick},
  lab/.style={font=\scriptsize, color=slate}
]
\node[d] (img) at (0,1.6) {\texttt{images/}\\$n$ 张};
\node[d] (gm)  at (0,0.5) {\texttt{gland\_masks/}\\同名};
\node[d] (em)  at (0,-0.6) {\texttt{eyelid\_masks/}\\同名};

\node[d] (idx) at (3.2,-0.05) {按 basename\\建立三元索引};
\node[d] (chk) at (6.3,-0.05) {完整性检查\\缺项/孤儿/重名/命名};
\node[bad] (fail) at (6.3,-1.8) {硬失败\\终止，不静默跳过};

\draw[arr] (img) -- (idx);
\draw[arr] (gm) -- (idx);
\draw[arr] (em) -- (idx);
\draw[arr] (idx) -- (chk);
\draw[arr] (chk) -- node[lab,right]{不通过} (fail);

\node[ok] (split) at (10.0,1.1) {患者级 holdout\\\nTrain/\nVal/\nTest\\seed 42};
\node[ok] (kf) at (10.0,-0.05) {患者分组\\复用同一编号};
\node[ok] (sig) at (10.0,-1.2) {文件清单签名 +\\派分签名（SHA-256）};
\draw[arr] (chk) -- node[lab,above]{通过} (split);
\draw[arr] (chk) -- (kf);
\draw[arr] (split) -- (sig);
\draw[arr] (kf) -- (sig);

\node[draw=alert!70, fill=alert!6, rounded corners=2pt, minimum width=0.95\linewidth, minimum height=7mm, align=center, font=\scriptsize] at (5.0,-3.0)
{签名随评估报告写入；缺失签名的报告判为无效，不得进入论文表格或权重转正};
\end{tikzpicture}}
\caption{三向配对、完整性与患者级划分。影像、腺体掩膜、睑缘掩膜按 basename 建立索引；完整性检查不通过时直接终止。划分采用患者级持久 holdout（\nTrain/\nVal/\nTest，\nTrainPatients/\nValPatients/\nTestPatients{} 名患者），同时记录文件清单签名与实际派分签名；5 折作为扩展复用同一患者分组。}
\label{fig:pairing}
\end{figure}


\subsection{分块尺寸的隐性影响}\label{sec:tiling}

\paragraph{实验口径。}本节的全部数值来自\textbf{同一权重、仅改变分块高度}的推理实验：阈值 0.7，TTA 关闭，不含后处理（不做 ROI 约束、不做连通域过滤、不做形态学平滑），在 40 例原始影像上评估。除分块高度外，预处理、权重、阈值与评价代码完全相同。这一口径是本节结论成立的前提；任何一项改变都不应与下表直接比较。

\begin{table}[t]
\centering
\caption{分块高度对 Dice 的影响（同一权重、仅改分块；阈值 0.7、TTA 关闭、无后处理与 ROI；$n=40$）。}
\label{tab:tiling}
\begin{tabular}{lccccccc}
\toprule
分块高度（像素） & 128 & 192 & 256 & 384 & 448 & 512 & 整图 $640\times1280$ \\
\midrule
Dice & 0.370 & 0.410 & 0.480 & \textbf{0.644} & 0.643 & 0.655 & 0.593 \\
\bottomrule
\end{tabular}
\end{table}

表~\ref{tab:tiling} 与图~\ref{fig:tiling} 给出七个真实测量点。三点值得强调。第一，从 128 到 384，Dice 提升了约 27 个百分点，这不是网络差异，而是同一权重的推理配置差异。第二，512 略高于 384，但整图推理反而降到 0.593：模型在 $384\times768$ 分块上训练，整图输入改变了有效感受野与上下文比例，训练—推理尺度不再匹配。第三，384 附近的平台（0.644/0.643/0.655）说明收益在某一区间饱和，继续增大分块并不单调变好。

\begin{figure}[t]
\centering
\begin{tikzpicture}
\begin{axis}[
  width=0.92\linewidth,
  height=62mm,
  xlabel={分块高度（像素）},
  ylabel={Dice（同权重，仅改分块）},
  xmin=0.5, xmax=7.5,
  ymin=0.33, ymax=0.69,
  xtick={1,2,3,4,5,6,7},
  xticklabels={128,192,256,384,448,512,整图},
  ytick={0.35,0.40,0.45,0.50,0.55,0.60,0.65},
  grid=major,
  grid style={draw=mist},
  axis line style={draw=slate},
  tick style={draw=slate},
  tick label style={font=\scriptsize},
  label style={font=\small},
  every axis plot/.append style={thick},
]
\addplot[color=teal, mark=*, mark size=2.2pt] coordinates {
  (1,\tileA) (2,\tileB) (3,\tileC) (4,\tileD) (5,\tileE) (6,\tileF) (7,\tileG)
};
\addplot[dashed, color=alert, domain=0.5:7.5] {\tileD};
\node[font=\scriptsize, color=alert, anchor=south west] at (axis cs:4.05,\tileD) {默认 384（0.644）};
\draw[-{Stealth[length=1.6mm]}, draw=slate] (axis cs:5.6,0.38) -- (axis cs:2.4,0.43);
\node[font=\scriptsize, color=slate, align=center] at (axis cs:5.4,0.375) {128$\to$384：+27 个百分点};
\end{axis}
\end{tikzpicture}
\caption{分块敏感性（同一权重、仅改变分块高度；阈值 0.7、TTA 关闭、不含后处理与 ROI；$n=\tilingN$）。384 为默认值；整图推理因训练—推理尺度不匹配反而下降。}
\label{fig:tiling}
\end{figure}


据此，本文把推理与训练的默认分块从早期的 192/256 统一修为 $384\times768$，并在所有报告中强制记录分块高度与重叠。该发现的方法学意义在于：若论文只报告“最优配置”而不报告分块，同一模型可以被叙述成强或弱，且读者无法复现；分块应被视为与主干、损失同等级的实验变量。

\subsection{Multi-Tile Ensemble：把分块从混杂变量变成推理工具}\label{sec:mte-system}

分块敏感性说明“单一分块”的结论不可迁移。与其只报告一个最优分块，我们把多个分块的概率图做平均：对 $h\in\{256,384,512\}$ 分别滑窗推理得到 $P_h$，MTE 输出 $\bar P=\frac{1}{3}\sum_h P_h$，再以阈值 0.7 二值化。MTE 不改变权重、不增加训练成本；推理成本约为单分块的 3 倍，低于 8 倍 TTA。它的作用不是“再涨一点 Dice”，而是把分块从不可控的混杂变量变成可显式控制的推理工具：当单分块 256 与 512 的结论冲突时，MTE 给出一个不依赖单点选择的稳定估计。

同时，三张 $P_h$ 的像素方差 $V=\mathrm{Var}_h(P_h)$ 是一个低成本的不确定性来源。我们把它与预测熵、8 倍 TTA 误差一起用于失败病例（Dice $<0.7$）检测，报告 AUROC。这样，分块既是被审计的对象，也是审计其他变量的工具。

\subsection{显存档位与消费级复现}\label{sec:vram}

为了让“可在消费级显卡上复现”成为可检验的主张，而不是口号，本文给出六档显存配方（表~\ref{tab:vram}、图~\ref{fig:vram}）。每档固定分块、batch、梯度累积、AMP、channels\_last、深监督与 TTA 开关，并在元数据中标注状态为\textbf{实测}或\textbf{估算}。实测值来自 RTX 4060 Laptop 8 GB、Python 3.11、torch 2.5.1 CUDA 12.1、700 训练样本（去重前）。当前去重后划分的训练集为 699 例，单例差异不影响吞吐结论。

\begin{table}[t]
\centering
\caption{显存档位配方。6 GB 与 8 GB 为实测，实测条件为去重前 700 训练样本；12 GB 及以上为外推估算，不得当作实测值引用。}
\label{tab:vram}
\small
\resizebox{\linewidth}{!}{%
\begin{tabular}{lcccccc}
\toprule
档位 & 分块 & batch/累积 & 深监督 & 峰值显存 & 吞吐（样本/秒） & 状态 \\
\midrule
6 GB  & $256\times512$  & 1/4 & 否 & 0.89 GB & 3.54 & 实测 \\
8 GB  & $384\times768$  & 1/4 & 是 & 1.227 GB & 3.10 & 实测 \\
12 GB & $384\times768$  & 2/2 & 是 & 估算 & 估算 & 估算 \\
16 GB & $512\times1024$  & 2/1 & 是 & 估算 & 估算 & 估算 \\
24 GB & $640\times1280$  & 2/1 & 是 & 估算 & 估算 & 估算 \\
40 GB+& $640\times1280$  & 4/1 & 是 & 估算 & 估算 & 估算 \\
\bottomrule
\end{tabular}}
\end{table}

8 GB 档的单轮实测为 225.6 s/epoch、3.10 样本/秒、峰值 1.227 GB；6 GB 档为 197.5 s/epoch、3.54 样本/秒、峰值 0.89 GB。上述 epoch 时间、吞吐与显存峰值均在同一批 700 个去重前训练样本上测得。对应实测 run 为 \path{20261106T203908Z_MGD-Former_nogit_7c6ef4}。两次冒烟均在预训练编码器未加载、随机初始化的条件下完成，日志记录了该降级原因。因此这些数字回答的是“能不能在 8 GB 上跑起来、吞吐多少”，不是“预训练后精度如何”。

\begin{figure}[t]
\centering
\resizebox{\linewidth}{!}{%
\begin{tikzpicture}[
  font=\scriptsize,
  meas/.style={draw=teal, fill=teal!25},
  est/.style={draw=slate!45, fill=mist, pattern=north east lines, pattern color=slate!35},
  lab/.style={anchor=west, font=\scriptsize},
  val/.style={anchor=west, font=\scriptsize, color=slate}
]
% axes
\draw[->, draw=slate] (0,0) -- (10.6,0) node[right, font=\scriptsize]{峰值显存（GB，示意比例）};
\foreach \y/\name/\peak/\sps/\stat in {
  0/6 GB/0.89/3.54/实测,
  1/8 GB/1.227/3.10/实测,
  2/12 GB/1.8/2.9/估算,
  3/16 GB/2.5/2.6/估算,
  4/24 GB/3.6/2.1/估算,
  5/40 GB+/5.0/1.8/估算} {
  \pgfmathsetmacro{\w}{(\peak)*3.2}
  \ifnum\y<2
    \draw[meas] (0,{0.55*\y+0.55}) rectangle (\w,{0.55*\y+1.05});
  \else
    \draw[est] (0,{0.55*\y+0.55}) rectangle (\w,{0.55*\y+1.05});
  \fi
  \node[anchor=east, font=\scriptsize] at (-0.1,{0.55*\y+0.8}) {\name};
  \ifnum\y<2
    \node[val] at (\w+0.1,{0.55*\y+0.8}) {\peak\,GB，\sps\,样本/秒（\stat）};
  \else
    \node[val] at (\w+0.1,{0.55*\y+0.8}) {\stat（不外推数值）};
  \fi
}
\node[anchor=west, font=\scriptsize, color=slate] at (0,-0.55) {实心：实测；斜线：估算。8 GB 实测 225.6\,s/epoch（去重前 700 训练样本，随机初始化）。};
\end{tikzpicture}}
\caption{六档显存配方。6 GB 与 8 GB 为 RTX 4060 Laptop 8 GB 上、去重前 700 训练样本口径的实测（峰值与吞吐），12 GB 及以上仅为档位规划，不外推具体数值。估算档不得在正文中冒充实测。}
\label{fig:vram}
\end{figure}


档位是安全边界而非建议值：选择档位会覆盖分块、batch 与 AMP 等字段，避免出现“能启动但必然 OOM”的组合。换机后首次训练必须观察 \texttt{peak\_vram\_gb}，实测后回填并改状态；估算档不得冒充实测。这条规则与第~\ref{sec:discipline}~节的报告纪律同源：可复现性依赖对“已知”与“外推”的显式区分。

\subsection{可编排流水线与实验登记}\label{sec:pipeline}

模型比较需要多折、多配置、可中断的长任务，手工记录极易丢失上下文。本文把训练—评估—统计—报告组织为可编排流水线，阶段为 \texttt{pretrain}、\texttt{train}、\texttt{finetune}、\texttt{eval}、\texttt{ablation}、\texttt{report}；阶段状态写入 \texttt{pipeline\_state.json}，中断后重跑自动跳过已完成条目，失败条目按配置重试。

每次运行生成独立目录，命名规则为
\begin{equation}
\texttt{run\_id}=\text{时间戳}\_\text{名称}\_\text{代码版本}\_\text{配置指纹前6位},
\end{equation}
重名自动追加序号，绝不覆盖已有 run。目录内固定包含：配置快照 \texttt{config.yaml}（含解析后的完整参数）、环境快照 \texttt{env.json}（代码版本、git commit、依赖、设备、路径）、逐轮指标 \texttt{metrics.jsonl}、完整训练状态 checkpoint、带 \texttt{split\_signature} 的 \texttt{reference\_report.json} 与结果摘要 \texttt{run\_result.json}（含峰值显存、吞吐、墙钟时间）。断点续训保存模型、优化器、调度器、AMP scaler、轮次、全局步数、随机数状态、EMA shadow 与数据签名。

实验登记的价值在于“多配置对比”可被机器读取。逐轮指标 \path{metrics.jsonl} 与结果摘要 \path{run_result.json} 既可支撑论文表格，也可支撑失败运行的追溯。第~\ref{sec:experiments}~节的实验矩阵直接映射到这些字段，而不是依赖人工笔记。

\begin{figure}[t]
\centering
\resizebox{\linewidth}{!}{%
\begin{tikzpicture}[
  p/.style={draw=slate, fill=mist, rounded corners=1.2pt, align=center, font=\scriptsize, minimum width=23mm, minimum height=9mm},
  s/.style={draw=teal, fill=teal!12, rounded corners=1.2pt, align=center, font=\scriptsize, minimum width=23mm, minimum height=9mm},
  r/.style={draw=alert!70, fill=alert!6, rounded corners=1.2pt, align=center, font=\scriptsize, minimum width=23mm, minimum height=9mm},
  arr/.style={-{Stealth[length=1.8mm]}, draw=slate, thick},
  lab/.style={font=\scriptsize, color=slate}
]
\node[p] (data) at (0,0) {\texttt{mgd-data check}\\完整性};
\node[p, right=5mm of data] (split) at (3.0,0) {患者级划分\\\nTrain/\nVal/\nTest};
\node[s, right=5mm of split] (sig) at (6.0,0) {文件清单 +\\派分签名};
\node[s, right=5mm of sig] (train) at (9.0,0) {每次 run 训练\\run 目录/EMA};
\node[s, right=5mm of train] (eval) at (12.0,0) {评估\\patient-level test};
\node[s, right=5mm of eval] (stat) at (15.0,0) {统计\\bootstrap/Wilcoxon};
\node[r, right=5mm of stat] (rep) at (18.0,0) {报告\\签名校验};
\draw[arr] (data)--(split);\draw[arr] (split)--(sig);\draw[arr] (sig)--(train);
\draw[arr] (train)--(eval);\draw[arr] (eval)--(stat);\draw[arr] (stat)--(rep);

\node[p, below=9mm of train] (state) {\texttt{pipeline\_state.json}\\阶段状态};
\node[p, below=9mm of eval] (retry) {失败重试\\断点续跑};
\draw[arr, dashed] (state) -- (train);
\draw[arr, dashed] (retry) -- (eval);

\node[lab, above=4mm of train] {配置快照 + env.json + metrics.jsonl};
\node[lab, above=4mm of eval] {分块 384×768、阈值 0.7、TTA 档位};
\node[lab, above=4mm of rep] {缺签名 $\Rightarrow$ 无效};
\end{tikzpicture}}
\caption{实验流水线。数据完整性检查与患者级划分先行，文件清单签名与派分签名随运行写入；每次训练产出独立 run 目录，评估与统计在 patient-level test 上进行；报告阶段校验签名，缺失即拒绝。状态文件支持失败重试与断点续跑。}
\label{fig:pipeline}
\end{figure}


\subsection{评测纪律与 \texttt{split\_signature}}\label{sec:discipline}

本文把评测写成协议，而不是“跑一次测试集”。协议包含五条：

\begin{enumerate}[label=\textbf{E\arabic*.}]
\item \textbf{患者级持久划分。}holdout 采用 \nTrain/\nVal/\nTest（\nTrainPatients/\nValPatients/\nTestPatients{} 名患者，seed 42），写入 \path{dataset_splits.json} 并同时记录文件清单签名与实际派分签名；5 折交叉验证复用同一份患者分组，不重新切分。
\item \textbf{签名强制。}\texttt{split\_signature} 是训练/验证/测试文件名有序集合的 SHA-256 摘要。任何评估报告必须携带该签名；缺失签名的报告判为无效，不得进入论文表格或模型转正流程。
\item \textbf{口径分离。}\texttt{train\_self} 必须显式标注，只能回答拟合程度；只有 holdout 或 5 折中未参与训练的折外预测可作为性能结论。
\item \textbf{统计完整。}均值给 bootstrap 1000 次重采样的 95\% 置信区间；模型间用逐样本配对 Wilcoxon 符号秩检验，并以 Holm 校正控制多重比较；报告有效样本数。
\item \textbf{权重状态。}训练产物一律标记 \texttt{verified=false}；只有带有效签名与独立评估的权重才能转正。加载未验证权重必须显式允许，并在日志与报告中写明原因。
\end{enumerate}

这条签名规则防的是三类真实风险：其一，调参后悄悄重切划分，使测试集事实上参与模型选择；其二，不同实验使用了不同的有效样本集合（因静默跳过损坏样本），却仍被放在同一张表中比较；其三，旧权重在全部数据上训练过，却被当作独立测试结果。签名把“哪一份样本参与了什么”变成不可篡改的元数据，从流程上排除这些风险。

\subsection{分层架构与部署形态}\label{sec:architecture}

系统按“内核—适配器”分层：\texttt{core} 不含任何 UI 依赖，包含配置、运行时、数据、模型、训练、推理、评估、成像、报告与工作流；\texttt{adapters} 提供 CLI 与 WebUI 薄壳，两者调用同一套 API。两端（Windows WebUI、Linux CLI）共享同一内核，版本 2.3.0，契约与回归测试共 146 项通过。

模型注册表统一管理 MGD-Former 与四条基线；输出契约由 \texttt{unwrap\_outputs()} 统一解包（旧模型返回张量，MGD-Former 返回 \texttt{logits} 与三个辅助头）。权重 manifest 记录文件名、架构、版本、SHA-256、输入尺寸、归一化口径（\texttt{scale01+clahe} 及 clip limit）、训练来源与 \texttt{verified} 状态；设备自适应覆盖 CUDA/CPU，推理服务对模型做单例缓存并支持显式释放。图~\ref{fig:system} 给出分层与数据流。

\begin{figure}[t]
\centering
\resizebox{\linewidth}{!}{%
\begin{tikzpicture}[
  layer/.style={draw=slate!60, fill=white, rounded corners=2pt, minimum width=150mm, minimum height=12mm, align=left, font=\scriptsize, inner sep=4pt},
  core/.style={draw=teal, fill=teal!8, rounded corners=2pt, minimum width=150mm, minimum height=16mm, align=left, font=\scriptsize, inner sep=4pt},
  arr/.style={-{Stealth[length=2mm]}, draw=slate, thick}
]
\node[layer] (ad) at (0,4.2) {\textbf{适配层}：CLI（\texttt{mgd-train} / \texttt{mgd-eval} / \texttt{mgd-pipeline} / \texttt{mgd-data} / \texttt{mgd-verify}）\quad WebUI（六分区）};
\node[core] (wf) at (0,2.4) {\textbf{workflow}：阶段编排（pretrain/train/finetune/eval/ablation/report）、多配置、多折、重试、断点续跑、\texttt{pipeline\_state.json}};
\node[core] (svc) at (0,0.7) {\textbf{core 服务}：training（specs/presets/engine/runner/EMA/losses）\quad inference（分块/TTA/阈值/ROI）\quad evaluation（Dice/IoU/HD95/ASSD/clDice/bootstrap/Wilcoxon）\quad reporting};
\node[layer] (pri) at (0,-0.9) {\textbf{原语}：data（pairing/splits/dataset/transforms/versioning）\quad models（MGD-Former / 四基线 / registry / manifest / verify）\quad imaging（preprocess / metrics / visualize）\quad runtime（device / errors / logging / doctor）};
\node[layer] (art) at (0,-2.3) {\textbf{资产}：\texttt{configs/}（YAML schema）\quad \texttt{runs/<run\_id>/}（配置快照、env、metrics、checkpoint、reference report）\quad 权重 manifest};
\draw[arr] (ad) -- (wf);
\draw[arr] (wf) -- (svc);
\draw[arr] (svc) -- (pri);
\draw[arr] (pri) -- (art);
\node[font=\scriptsize, color=slate, anchor=west] at (-7.7,5.0) {设备自适应：CUDA / CPU；模型单例缓存与显式释放；未验证权重必须显式允许并标注};
\end{tikzpicture}}
\caption{系统分层与数据流。适配层只做薄壳，所有入口经过同一训练/推理/评估实现；workflow 负责编排与续跑；资产层保存配置快照、逐轮指标、完整 checkpoint 与带签名的评估报告。}
\label{fig:system}
\end{figure}


这些实现细节之所以进入方法章，是因为它们决定了实验能否被他人重跑：同一份 manifest、同一份配置快照、同一套分块默认值与同一份签名，构成“结果—代码—数据”的可追溯链条。缺少任何一环，第~\ref{sec:results}~节的比较都只能是不可复核的数字。
